\documentclass[10pt,a4paper]{article}

\usepackage[a4paper,top=0.82cm,bottom=0.82cm,left=1.05cm,right=1.05cm]{geometry}
\usepackage{fontspec}
\usepackage[french]{babel}
\usepackage[hidelinks]{hyperref}
\usepackage{xcolor}
\usepackage{titlesec}
\usepackage{enumitem}

\setmainfont{Arial}
\pagestyle{empty}
\setlength{\parindent}{0pt}
\setlength{\parskip}{0pt}
\linespread{0.95}
\hyphenpenalty=10000
\exhyphenpenalty=10000
\sloppy
\setlist[itemize]{leftmargin=1.05em,itemsep=0pt,topsep=1pt,parsep=0pt,label=-}
\urlstyle{same}

\definecolor{heading}{HTML}{123867}
\titleformat{\section}{\large\bfseries\color{heading}}{}{0pt}{}[\vspace{-5pt}\rule{\linewidth}{0.5pt}]
\titlespacing{\section}{0pt}{5pt}{3pt}

\newcommand{\role}[3]{%
  \textbf{#1} \hfill #3\\
  \textit{#2}\vspace{1pt}
}

\newcommand{\edu}[4]{%
  \textbf{#1} -- #2 \hfill #3 #4\\[5pt]
}

\begin{document}

\begin{center}
{\LARGE \textbf{Kossi Richard Allado}}\\[2pt]
{\large \textbf{Stage PFA / Observation | Support IT, Réseaux \& Systèmes}}\\[4pt]
Beni Mellal / Casablanca \quad | \quad +212 621 074 305 \quad | \quad \href{mailto:alladokossirichard2025@gmail.com}{alladokossirichard2025@gmail.com}\\[2pt]
\href{https://allakori.github.io/}{allakori.github.io} \quad | \quad \href{https://github.com/ALLAKORI}{github.com/ALLAKORI} \quad | \quad \href{https://www.linkedin.com/in/kossi-richard-allado-50177b2a5}{linkedin.com/in/kossi-richard-allado}\\[3pt]
\textit{Disponible à partir du lundi 8 juin 2026 -- Stage de 1 à 2 mois}
\end{center}
\vspace{-3pt}

\section{Profil}
Étudiant ingénieur en cybersécurité et intelligence artificielle à l'ENSA Béni Mellal, avec de bonnes bases en support informatique, réseaux, systèmes et diagnostic technique. Intéressé par un stage PFA / observation chez Intelcia pour assister les techniciens IT, participer à l'installation/configuration des postes, contribuer au traitement des incidents de premier niveau et documenter les interventions.

\section{Compétences techniques}
\textbf{Support IT :} assistance utilisateur, troubleshooting, suivi d'incidents, documentation, installation/configuration.\\
\textbf{Réseaux :} adressage IP, DNS/DHCP, SSH, analyse de trafic, Nmap, Wireshark, diagnostic de connectivité.\\
\textbf{Systèmes :} Linux, notions Windows, comptes utilisateurs, services, journaux système, hardening de base.\\
\textbf{Outils / automatisation :} Bash, Python, SQL, Git/GitHub, rapports techniques, tests et validation.

\section{Projets techniques}
\textbf{Linux Security Labs : administration, audit et diagnostic système} \hfill Linux / Systèmes\\
\href{https://github.com/ALLAKORI/linux-security-labs}{github.com/ALLAKORI/linux-security-labs}
\begin{itemize}
  \item Mis en pratique des tâches Linux : configuration, audit système, analyse de logs et vérification de services.
  \item Utilisé auditd, Lynis, commandes Linux et documentation technique pour détecter des problèmes et proposer des corrections.
\end{itemize}

\textbf{Cybersecurity \& IAM Labs} \hfill Comptes / Accès / Sécurité\\
\href{https://github.com/ALLAKORI/cybersecurity-iam-labs}{github.com/ALLAKORI/cybersecurity-iam-labs}
\begin{itemize}
  \item Déployé un environnement IAM avec Keycloak : utilisateurs, rôles, MFA/TOTP, RBAC et tests de contrôle d'accès.
  \item Manipulé SSH, OpenSSL/PKI, notions de sécurité réseau et tests d'accès dans un environnement de laboratoire.
\end{itemize}

\textbf{Investigation de logs et diagnostic d'incident} \hfill Logs / Support technique\\
\href{https://allakori.github.io/writeups/tryhackme/volt-typhoon/}{allakori.github.io/writeups/tryhackme/volt-typhoon}
\begin{itemize}
  \item Analysé des journaux Windows/applicatifs avec Splunk pour reconstruire une chronologie d'incident.
  \item Produit un rapport clair avec événements importants, indicateurs, hypothèses et actions de remédiation.
\end{itemize}

\section{Expérience}
\role{Stagiaire IA / Maintenance prédictive}{ONEE - Branche Eau, Kasba Tadla}{Août 2025}
\begin{itemize}
  \item Préparé et analysé des données capteurs de pompe haute pression pour prédire des défauts potentiels.
  \item Construit un prototype Python avec pipeline de préparation de données, modèle Random Forest et interface de test.
  \item Développé une méthode de diagnostic et de présentation des résultats exploitable par un utilisateur technique.
\end{itemize}

\section{Formation}
\edu{Cycle ingénieur}{Cybersécurité et Intelligence Artificielle, ENSA Béni Mellal}{2024 -- Présent}{}
\edu{DEUST}{Mathématiques, Informatique, Physique, FST Fès}{2022 -- 2024}{-- Mention Bien}
\edu{Baccalauréat scientifique}{Sciences expérimentales, Lycée de Djagblé, Togo}{2022}{-- Mention Très Bien}

\section{Certifications}
TryHackMe Junior Penetration Tester Learning Path (en cours) | Udemy Linux Command Line | Google Cybersecurity: Foundations | Cisco Ethical Hacker

\section{Langues}
Français : courant | Anglais : professionnel technique | Éwé : langue maternelle

\end{document}
